% =============================================================================
% BLOCO 7 — SÍNTESE E ENCERRAMENTO  (17:50–18:00)
% =============================================================================
\secao{Encerramento}{Síntese: qual técnica usar?}{17:50 -- 18:00}

\begin{frame}{As quatro técnicas lado a lado}
  \scriptsize
  \renewcommand{\arraystretch}{1.35}
  \newcolumntype{Y}{>{\raggedright\arraybackslash}X}
  \begin{tabularx}{\textwidth}{@{}>{\raggedright\arraybackslash}p{1.9cm}YYYY@{}}
    \toprule
    & \textbf{Experimento} & \textbf{Pareamento (PSM)} & \textbf{Dif.-em-dif.} & \textbf{Regr. descontínua} \\ \midrule
    \textbf{Contrafactual vem de\ldots} & Sorteio entre elegíveis & Não participantes \emph{parecidos} & Tendência do grupo de controle & Quem está do outro lado do corte \\
    \textbf{Pressuposto-chave} & Sorteio bem executado & Sem confundidores \emph{não observados}; região comum & \textbf{Tendências paralelas} & Sem manipulação do corte; continuidade \\
    \textbf{Dados necessários} & Baseline e follow-up desenhados \emph{ex-ante} & Muitas covariáveis; um período & Pelo menos 2 períodos, tratados e controles & Índice contínuo e corte claro; muitas obs. no corte \\
    \textbf{Teste/diagnóstico} & Balanceamento; \emph{compliance}; atrito & Balanceamento (SMD); região comum & Pré-tendências; placebo & Janela; covariáveis; manipulação; placebo \\
    \textbf{Limitação} & Custo, ética, validade externa & Viés se faltar variável relevante & Choques diferenciais & Efeito \emph{local} \\
    \textbf{Em nossa simulação} & R\$ 295--302 & R\$ 297 (sem baseline: 476) & R\$ 300 (violado: 420) & 6,2 p.p. (verdade: 6) \\ \bottomrule
  \end{tabularx}
  \fonte{Batista e Domingos (2017), pp.~6--13; simulações dos scripts 02 a 05.}
\end{frame}

\begin{frame}{Qual técnica usar? Um roteiro de decisão}
  \centering
  \begin{tikzpicture}[node distance=0.55cm,scale=0.82,every node/.style={transform shape},
      dec/.style={draw=laranja,fill=laranjaclaro,rounded corners=4pt,thick,align=center,inner sep=5pt,font=\small,minimum width=4.2cm},
      res/.style={draw=indigo,fill=indigoclaro,rounded corners=4pt,thick,align=center,inner sep=5pt,font=\small\bfseries,minimum width=3.8cm}]
    \node[dec] (d1) at (0,3.6) {A avaliação está sendo \textbf{planejada} antes\\da política e é possível \textbf{sortear}?};
    \node[res] (r1) at (7.5,3.6) {Experimento (RCT)};
    \node[dec] (d2) at (0,1.8) {Há um \textbf{corte} de elegibilidade\\em um índice contínuo?};
    \node[res] (r2) at (7.5,1.8) {Regressão descontínua};
    \node[dec] (d3) at (0,0) {Há dados \textbf{antes e depois},\\para tratados e não tratados?};
    \node[res] (r3) at (7.5,0) {Diferença-em-diferenças};
    \node[dec] (d4) at (0,-1.8) {Observamos bem os fatores\\que explicam a participação?};
    \node[res] (r4) at (7.5,-1.8) {Pareamento / PSM\\\scriptsize (ou DD com pareamento)};
    \draw[setal] (d1) -- node[above,font=\scriptsize]{sim} (r1);
    \draw[setal] (d2) -- node[above,font=\scriptsize]{sim} (r2);
    \draw[setal] (d3) -- node[above,font=\scriptsize]{sim} (r3);
    \draw[setal] (d4) -- node[above,font=\scriptsize]{sim} (r4);
    \draw[seta] (d1) -- node[right,font=\scriptsize]{não} (d2);
    \draw[seta] (d2) -- node[right,font=\scriptsize]{não} (d3);
    \draw[seta] (d3) -- node[right,font=\scriptsize]{não} (d4);
    \node[font=\small\itshape,text=cinza,below=0.15cm of d4,align=center] {não $\Rightarrow$ descrição, qualitativo, resultados\\(sem pretensão causal)};
  \end{tikzpicture}
  \fonte{Síntese nossa a partir de Batista e Domingos (2017). Na prática, combine sempre com métodos qualitativos.}
\end{frame}

\begin{frame}{Cinco perguntas antes de confiar em uma estimativa}
  \begin{enumerate}
    \item \textbf{Qual é o contrafactual?} De onde ele vem e por que é crível?
    \item \textbf{Quais grupos estão sendo comparados} e eles eram parecidos \emph{antes}? (balanceamento, pré-tendências)
    \item \textbf{Qual pressuposto sustenta o resultado?} Ele foi testado ou só declarado?
    \item \textbf{Para quem vale o efeito?} (população do experimento, tratados, \emph{no corte})
    \item \textbf{O que o número não mostra?} Quem ganhou mais, quem ganhou menos, e \emph{por quê}?
  \end{enumerate}
  \vspace{0.3cm}
  \begin{block}{Lembrete do artigo}
    \small A discricionariedade do avaliador é grande: o mesmo programa pode ser recomendado ou interrompido dependendo da técnica escolhida (vimos 45 $\times$ 5 $\times$ 15 no DiD).
  \end{block}
\end{frame}

\begin{frame}{Mensagem final: mais que boas intenções}
  \begin{columns}[T]
    \begin{column}{0.58\textwidth}
      \begin{itemize}
        \item O experimento é o \emph{padrão ouro}, mas nem sempre viável: por isso existem as técnicas \ind{quase-experimentais};
        \item Elas \textbf{emulam} um grupo de controle, cada uma com seus pressupostos: \destaque{diga quais são, teste o que puder};
        \item Estimar o impacto \ind{médio} é essencial para decidir manter ou interromper políticas --- mas ele não mostra as \emph{pessoas por trás dos números};
        \item Combinar métodos aumenta a contribuição de cada um: \textbf{quanto} e \textbf{por quê}.
      \end{itemize}
    \end{column}
    \begin{column}{0.42\textwidth}
      \begin{beamercolorbox}[wd=\textwidth,rounded=true,sep=0.35cm]{laranjabox}
        \small\textbf{Para levar para casa}
        \begin{itemize}
          \item Rodar de novo os scripts, mudando parâmetros;
          \item Escolher uma política que vocês estudam e responder às 5 perguntas;
          \item Desenhar, ainda que em esboço, como \emph{seria} a avaliação ideal.
        \end{itemize}
      \end{beamercolorbox}
    \end{column}
  \end{columns}
  \fonte{Batista e Domingos (2017), Conclusão (p.~20).}
\end{frame}

\begin{frame}{Para continuar: leituras}
  \small
  \begin{itemize}
    \item \textbf{Batista, M. e Domingos, A. (2017).} Mais que boas intenções: técnicas quantitativas e qualitativas na avaliação de impacto de políticas públicas. \emph{Revista Brasileira de Ciências Sociais}, 32(94), e329414. \cod{doi.org/10.17666/329414/2017}
    \item \textbf{Gertler, P. et al. (2011).} \emph{Impact Evaluation in Practice}. World Bank.
    \item \textbf{Khandker, S., Koolwal, G. e Samad, H. (2010).} \emph{Handbook on Impact Evaluation: Quantitative Methods and Practices}. World Bank.
    \item \textbf{Glennerster, R. e Takavarasha, K. (2013).} \emph{Running Randomized Evaluations}. Princeton UP.
    \item \textbf{Dunning, T. (2012).} \emph{Natural Experiments in the Social Sciences}. Cambridge UP.
    \item \textbf{Ipea / Casa Civil.} \emph{Avaliação de Políticas Públicas: Guia Prático de Análise ex-Ante} e \emph{ex-Post} (gratuitos, no portal do governo federal).
    \item \textbf{Patton, M. (2002).} \emph{Qualitative Research and Evaluation Methods}. Sage.
    \item \textbf{Bauer, M. e Gaskell, G. (2010).} \emph{Pesquisa qualitativa com texto, imagem e som}. Vozes.
  \end{itemize}
\end{frame}

\begin{frame}{Repositório do curso e contato}
  \begin{columns}[T]
    \begin{column}{0.55\textwidth}
      \begin{itemize}
        \item \textbf{GitHub}:\ {\footnotesize\cod{github.com/andersonheri/avaliacao-pp-each}}
        \item Pastas: \cod{slides/}, \cod{scripts/} (R comentado), \cod{atividades/}, \cod{gabaritos/}, \cod{material/}
        \item Tudo é reproduzível: mudem as sementes, os efeitos, os tamanhos de amostra
      \end{itemize}
      \vspace{0.3cm}
      \textbf{Anderson Henrique}\\
      \small Pós-doutorando CEM/USP (FAPESP/CEPid) \,$\cdot$\, INCT-QualiGov \,$\cdot$\, IPEA\\
      \cod{andersonheri@gmail.com}
    \end{column}
    \begin{column}{0.45\textwidth}
      \begin{atividadebox}{Feedback (2 min)}
        \small
        O que funcionou? O que mudariam? Que técnica vocês querem aprofundar?\\[0.2cm]
        Respondam no chat ou por e-mail.
      \end{atividadebox}
    \end{column}
  \end{columns}
\end{frame}

{\setbeamercolor{background canvas}{bg=indigo}
\begin{frame}[plain,noframenumbering]
  \begin{tikzpicture}[remember picture,overlay]
    \fill[laranja] (current page.south west) rectangle ++(0.5,\paperheight);
    \node[anchor=west,text=white,font=\fontsize{34}{38}\selectfont\bfseries,align=left] at ([xshift=1.5cm,yshift=0.5cm]current page.west) {Obrigado!};
    \node[anchor=west,text=laranjaclaro,font=\Large,align=left] at ([xshift=1.5cm,yshift=-1cm]current page.west) {Perguntas, críticas e sugestões};
  \end{tikzpicture}
\end{frame}}
